\documentclass[preview]{standalone}

\usepackage{amsmath}
\usepackage{amssymb}
\usepackage{cancel}
\usepackage{stellar}
\usepackage{definitions}
\usepackage{bettelini}

\begin{document}

\id{surface-flux}
\genpage

\section{Flux Through a Surface}

\begin{snippetdefinition}{surface-flux-definition}{Flux through an oriented surface}
    Let \(\Sigma\subseteq\realnumbers^3\) be an oriented regular surface
    with unit normal \(\nu\), and let
    \(F\colon U\subseteq\realnumbers^3\fromto\realnumbers^3\) be a
    continuous vector field on a neighborhood of \(\Sigma\). The
    \emph{flux of \(F\) through \(\Sigma\)} is
    \[
        \iint_\Sigma F\cdot\nu\,dS.
    \]
\end{snippetdefinition}

\begin{snippet}{surface-flux-parametrization-formula}
    If the chosen orientation is induced by a parametrization
    \(\sigma\colon\overline A\fromto\Sigma\), then
    \begin{align*}
        \iint_\Sigma F\cdot\nu\,dS
        &=\iint_A F(\sigma(u,v))\cdot
        \frac{\sigma_u\times\sigma_v}
        {\cancel{\|\sigma_u\times\sigma_v\|}}
        \cancel{\|\sigma_u\times\sigma_v\|}\,du\,dv\\
        &=\iint_A F(\sigma(u,v))\cdot
        (\sigma_u\times\sigma_v)\,du\,dv.
    \end{align*}
    Reversing the orientation changes the sign of the flux.
\end{snippet}

\begin{snippet}{surface-flux-as-two-form}
    For \(F=(F_1,F_2,F_3)\), the scalar triple product expands as the
    pullback of the \(2\)-form
    \[
        \omega_F=
        F_3\,dx\wedge dy-F_2\,dx\wedge dz+F_1\,dy\wedge dz.
    \]
    Indeed,
    \[
        \omega_F(\sigma_u,\sigma_v)
        =F(\sigma)\cdot(\sigma_u\times\sigma_v),
    \]
    so the flux integral is equivalently \(\iint_\Sigma\omega_F\).
\end{snippet}

\section{Examples}

\begin{snippetexample}{vertical-field-sphere-flux-example}{Flux through a sphere}
    Let \(F(x,y,z)=(0,0,z)\), and orient the sphere \(S_R\) outward. With
    spherical coordinates,
    \[
        \sigma_\varphi\times\sigma_\theta
        =R\sin\varphi\,\sigma(\varphi,\theta).
    \]
    Therefore
    \begin{align*}
        \iint_{S_R}F\cdot\nu\,dS
        &=\integral[0][\pi]
        [\integral[0][2\pi]
        [R^3\sin\varphi\cos^2\varphi][\theta]][\varphi]\\
        &=\frac43\pi R^3.
    \end{align*}
    This equals the volume of the ball because
    \(\operatorname{div}F=1\), as predicted by the divergence theorem.
\end{snippetexample}

\begin{snippetexample}{cylinder-outward-flux-example}{Outward flux through a cylinder}
    Consider the lateral surface
    \[
        x^2+y^2=R^2,
        \qquad 0\leq z\leq1,
    \]
    oriented radially outward, and the field
    \(F(x,y,z)=(\sin z,y,x^2)\). With
    \(\sigma(\theta,z)=(R\cos\theta,R\sin\theta,z)\),
    \[
        \sigma_\theta\times\sigma_z
        =(R\cos\theta,R\sin\theta,0).
    \]
    Hence
    \begin{align*}
        \iint_\Sigma F\cdot\nu\,dS
        &=\integral[0][2\pi]
        [\integral[0][1]
        [R\sin z\cos\theta+R^2\sin^2\theta][z]][\theta]\\
        &=\pi R^2.
    \end{align*}
    On the lower and upper disks, the outward normals are respectively
    \((0,0,-1)\) and \((0,0,1)\). Their fluxes are
    \[
        -\iint_{x^2+y^2\leq R^2}x^2\,dx\,dy
        =-\frac{\pi R^4}{4},
        \qquad
        \iint_{x^2+y^2\leq R^2}x^2\,dx\,dy
        =\frac{\pi R^4}{4},
    \]
    so they cancel. The total outward flux through the closed cylinder is
    \(\pi R^2\), its volume, since \(\operatorname{div}F=1\).
\end{snippetexample}

\end{document}
